\subsection{Modelagem Computacional da Evolução da Capacidade Resistente de Vigas de Concreto Armado sob Corrosão}

Este trabalho implementa um modelo numérico para avaliar, ao longo do tempo, a degradação da capacidade resistente à flexão de vigas de concreto armado submetidas à corrosão das armaduras. A formulação combina (i) os modelos normativos de cálculo da resistência à flexão conforme a NBR 6118:2023 e (ii) a modelagem físico-química dos efeitos da corrosão, com base em Azad e Al-Gohi (2008) e Peng e Stewart (2016). 

Apesar de o algoritmo completo realizar múltiplas simulações estocásticas em diversos passos de tempo, apresenta-se aqui apenas um passo de tempo e uma única simulação, a fim de descrever de forma clara o procedimento computacional.

\subsection{Parâmetros iniciais e propriedades mecânicas}

A viga analisada apresenta os seguintes parâmetros geométricos e mecânicos:

\begin{itemize}
    \item largura da seção: $b_w$;
    \item altura total: $h$;
    \item profundidade útil: $d = \alpha h$, com $\alpha = 0{,}90$;
    \item cobrimento nominal: $c_{\mathrm{nom}} = 25\,\text{mm}$;
    \item resistência característica do concreto: $f_{ck}$;
    \item resistência característica do aço: $f_{yk} = 500\,\text{MPa}$;
    \item módulo de elasticidade do aço: $E_s = 200\,\text{GPa}$;
    \item momento característico das ações permanentes: $M_{Gk} = 6{,}32\,\text{kN·m}$;
    \item momento característico das ações variáveis: $M_{Qk} = 1{,}381\,\text{kN·m}$.
\end{itemize}

A área inicial de armadura tracionada é

\[
A_{s,0} = 1{,}6 \times 10^{-4}\,\text{m}^2.
\]

Considerando um diâmetro nominal das barras igual a $d_b = 12{,}5\,\text{mm}$, obtém-se o número de barras:

\[
n_b = \frac{A_{s,0}}{\dfrac{\pi d_b^2}{4}}.
\]

A progressão da frente de carbonatação é modelada por:

\[
y_{\mathrm{carb}}(t) = k_c \sqrt{t}, 
\qquad
k_c = 6\,\text{mm}/\sqrt{\text{ano}}.
\]

\subsection{Modelo mecânico segundo a NBR 6118}

\subsubsection{Equilíbrio de forças na seção}

No cálculo da resistência última à flexão utiliza-se o equilíbrio entre as forças resistentes de compressão no concreto e de tração no aço:

\[
R_c(\beta) - R_s(\beta) = 0,
\]

em que $\beta = x/d$ representa a profundidade relativa da linha neutra. A força de compressão no concreto é expressa como

\[
R_c = \sigma_{cd}\, \lambda_c x b_w,
\]

onde $\sigma_{cd}$ é a tensão de cálculo do concreto comprimido. A força resistente no aço é dada por

\[
R_s = \sigma_s A_s,
\]

onde $\sigma_s$ é a tensão no aço obtida a partir do diagrama elasto-plástico definido pela norma.

A profundidade da linha neutra $x$ resulta da solução numérica da equação de equilíbrio:

\[
f(\beta) = R_c(\beta) - R_s(\beta) = 0.
\]

\subsubsection{Momento resistente sem corrosão}

Uma vez determinada a profundidade da linha neutra, o momento resistente inicial é:

\[
M_{Rd,0} = R_c \left( d - \frac{1}{2}\lambda_c x \right).
\]

\subsection{Modelagem da corrosão}

Após o início da corrosão, considerado no tempo $t_{\mathrm{ini}}$, o modelo incorpora efeitos de (i) perda de seção das barras e (ii) redução da aderência aço-concreto.

\subsubsection{Índice de corrosão ajustado pela temperatura}

Segundo Peng e Stewart (2016), o índice de corrosão ajustado à temperatura ambiente é:

\[
i_{\mathrm{corr}}(T) = i_{\mathrm{corr},20} \left[ 1 + k(T - 20) \right],
\]

com

\[
k =
\begin{cases}
0{,}073, & T > 20^\circ\text{C}, \\
0{,}025, & T < 20^\circ\text{C}, \\
0, & T = 20^\circ\text{C}.
\end{cases}
\]

\subsubsection{Perda de diâmetro da barra}

Definindo o período de corrosão como

\[
t_{\mathrm{corr}} = t - t_{\mathrm{ini}},
\]

a redução do diâmetro original $d_0$ (em mm) é dada por (Azad \& Al-Gohi, 2008):

\[
\Delta d = 0{,}0232 \, i_{\mathrm{corr}} \, t_{\mathrm{corr}}.
\]

O diâmetro residual da barra torna-se:

\[
d(t) = d_0 - \Delta d.
\]

A área efetiva da armadura é então:

\[
A_s(t) = n_b \, \frac{\pi d(t)^2}{4}.
\]

\subsubsection{Redução da aderência no aço}

O coeficiente de perda de aderência utilizado no modelo é:

\[
c_f =
\min\left[
1,\,
\dfrac{5}{\, d_0^{0{,}54} \left( i_{\mathrm{corr}}^\ast \, t_{\mathrm{corr}} \, 365 \right)^{0{,}19}}
\right],
\]

onde $i_{\mathrm{corr}}^\ast = i_{\mathrm{corr}}/1000$ é expresso em A/mm\(^2\).

\subsubsection{Momento resistente com corrosão}

A resistência à flexão degradada é calculada como:

\[
M_{Rd,\mathrm{cor}}(t) = c_f \, M_{Rd}(A_s(t)).
\]

\subsubsection{Cálculo da margem de segurança}

Uma vez obtido o momento resistente degradado $M_{Rd,\mathrm{cor}}(t)$, avalia-se o estado limite último de flexão mediante a comparação entre resistência e ações solicitantes. O momento solicitante total é definido como

\[
M_S = M_{Gk} + M_{Qk},
\]

onde $M_{Gk}$ e $M_{Qk}$ são, respectivamente, os momentos característicos das ações permanentes e variáveis.

A margem de segurança no instante $t$ é dada por

\[
G(t) = M_{Rd,\mathrm{cor}}(t) - M_S.
\]

Trata-se de uma função estado, amplamente utilizada em análises de confiabilidade estrutural. A interpretação física é direta:

\[
\begin{cases}
G(t) \ge 0, & \text{estrutura segura (estado não violado)}, \\[4pt]
G(t) < 0,  & \text{ocorrência do estado limite último}.
\end{cases}
\]

Assim, o valor de $G(t)$ indica a distância, em termos de momento fletor resistente, entre a capacidade residual da seção corroída e a demanda estrutural. A progressiva redução de $G(t)$ ao longo do tempo reflete a deterioração da barra longitudinal, tanto pela perda de seção quanto pela redução da aderência aço–concreto.

Este procedimento é repetido para todos os tempos da simulação e para todas as amostras estocásticas, permitindo a determinação de curvas de confiabilidade e da probabilidade de falha em cada ano de vida útil.

\subsection{Exemplo numérico para um único passo de tempo}

Considere o passo de tempo $t = 40$ anos, com temperatura ambiente $T = 23{,}56^\circ$C, índice de corrosão a 20°C igual a $i_{\mathrm{corr},20} = 0{,}431\,\mu\text{A/cm}^2$, diâmetro original $d_0 = 12{,}5\,\text{mm}$ e início da corrosão em $t_{\mathrm{ini}} = 25$ anos.

\subsubsection{Verificação da carbonatação}

\[
y_{\mathrm{carb}}(40) = 6\sqrt{40} = 37{,}95\,\text{mm}.
\]

Como $y_{\mathrm{carb}}(40) > c_{\mathrm{nom}} = 25\,\text{mm}$, considera-se que já ocorreu o início da corrosão.

\subsubsection{Índice de corrosão ajustado}

\[
i_{\mathrm{corr}}(23{,}56)
= 0{,}431 \, \left[ 1 + 0{,}073 (23{,}56 - 20) \right]
= 0{,}542\,\mu\text{A/cm}^2.
\]

\subsubsection{Perda de diâmetro}

\[
t_{\mathrm{corr}} = 40 - 25 = 15\,\text{anos},
\]

\[
\Delta d = 0{,}0232 \times 0{,}542 \times 15 = 0{,}1886\,\text{mm},
\]

\[
d(40) = 12{,}5 - 0{,}1886 = 12{,}3114\,\text{mm}.
\]

\subsubsection{Cálculo da área efetiva de aço}

A área da barra corroída é:

\[
A_s(40)
= \frac{\pi}{4}\,d(40)^2
= \frac{\pi}{4}(12{,}3114^2)
= 119{,}06\,\text{mm}^2.
\]

A barra original ($d_0 = 12{,}5$ mm) possuía:

\[
A_{s,0} = 122{,}72\,\text{mm}^2.
\]

Logo, há uma redução de:

\[
\Delta A_s = A_{s,0} - A_s(40) = 3{,}66\,\text{mm}^2.
\]

\subsubsection{Cálculo do momento resistente}

Com a perda de área, o momento resistente reduzido é obtido pelo equilíbrio seccional:

\[
M_{Rd}(40)
= 0{,}85\, f_{cd} \, b_w \, x \left( d - \frac{x}{2} \right),
\]

com reposicionamento neutro recalculado usando $A_s(40)$.

Considerando os parâmetros usados no código (que retornam o valor numérico):

\[
M_{Rd}(40) = 34{,}82\,\text{kN·m}.
\]

\subsubsection{Coeficiente de aderência}

O coeficiente redutor é:

\[
c_f = 
\dfrac{5}
{ d_0^{0{,}54}
\left[
(i_{\mathrm{corr}}/1000)\, t_{\mathrm{corr}}\, 365
\right]^{0{,}19}}
=
\dfrac{5}
{ 12{,}5^{0{,}54}
\left[
(0{,}542/1000)\cdot 15 \cdot 365
\right]^{0{,}19}}
= 0{,}873.
\]

Assim:

\[
M_{Rd,\mathrm{cor}}(40)
= c_f \, M_{Rd}(40)
= 0{,}873 \times 34{,}82
= 30{,}39\,\text{kN·m}.
\]

\subsection{Margem de segurança}

Os momentos solicitantes adotados no exemplo (conforme o código) são:

\[
M_{Gk} = 12{,}0\,\text{kN·m}, \qquad
M_{Qk} = 8{,}0\,\text{kN·m}.
\]

Logo:

\[
M_{Gk} + M_{Qk} = 20{,}0\,\text{kN·m}.
\]

A margem de segurança é então:

\[
G(40)
= M_{Rd,\mathrm{cor}}(40) - \bigl( M_{Gk} + M_{Qk} \bigr)
= 30{,}39 - 20{,}0
= 10{,}39\,\text{kN·m}.
\]

\[
\boxed{G(40) = 10{,}39\,\text{kN·m}}
\]

Esse resultado indica que, aos 40 anos, mesmo considerando corrosão, a seção ainda possui capacidade resistente superior às solicitações atuantes.
